\documentclass[11pt,a4paper]{article}

\usepackage[utf8]{inputenc}
\usepackage[T1]{fontenc}
\usepackage[english]{babel}
\usepackage{amsmath,amssymb,amsfonts}
\usepackage{graphicx}
\usepackage{booktabs}
\usepackage{array}
\usepackage{float}
\usepackage[margin=2.2cm]{geometry}
\usepackage[colorlinks=true,linkcolor=blue,citecolor=blue,urlcolor=blue]{hyperref}

\graphicspath{{./figures/}}

\title{\textbf{Nodal Vector Meshless Method for Electromagnetic Eigenvalue Problems in Discontinuous and Anisotropic Media}}

\author{
  Luilly A. G. Ortiz \and Renato C. Mesquita \and Barbara M. F. Gon\c{c}alves \\
  \small Universidade Federal de Minas Gerais (UFMG), School of Engineering \\
  \small Graduate Program in Electrical Engineering, Belo Horizonte, MG, Brazil
}

\date{}

\begin{document}

\maketitle

\begin{abstract}
This work presents the formulation and systematic validation of the \textbf{2D Vector Nodal Meshless Method (VNMM)} for electromagnetic eigenvalue problems in heterogeneous, discontinuous, and anisotropic media on unstructured point clouds. Each node carries a single scalar degree of freedom representing the electric field projection along a unit director vector ($\mathbf{t}_i$), combining a complete linear polynomial vector basis ($\mathcal{P}^1$), $k$-d tree local support selection, and variational div-curl regularization. To handle material interfaces without numerical smearing or Gibbs oscillations, a \textbf{strict material support partition rule} is introduced: stencils are confined to their own material subdomain, and interface nodes are assigned purely tangential directors ($\mathbf{t}_j = \boldsymbol{\tau}$), enforcing exact $C^0$ continuity in $E_t$ while allowing sharp normal jumps ($[E_n] \neq 0$). Benchmarks verify machine-precision ($10^{-16}$) Patch Test completeness, optimal $\mathcal{O}(h^2)$ convergence, over 50\% mean eigenvalue error reduction on moderately irregular clouds (10\%--25\% jitter) compared to regular grids, and spurious-free accuracy across wide permittivity contrasts ($\varepsilon_r \in [2.25, 16.0]$), gyrotropic H-shaped waveguides, curved PEC inclusions, and a $9:1$ four-material corner junction ($0.34\%$ dominant-mode error).
\end{abstract}

\noindent\textbf{Keywords:} Vector Nodal Meshless Method, Point Clouds, $H(\mathbf{curl})$ Spaces, Anisotropic Media, Multimaterial Interfaces, Electromagnetic Eigenvalue Problems, Spurious-Free Formulations.

\section{Introduction}

The numerical simulation of high-frequency electromagnetic devices---such as dielectric-loaded waveguides, resonant cavities, and gyrotropic components---relies on weak formulations of Maxwell's equations posed in the Sobolev space $H(\mathbf{curl})$ \cite{harrington2001,collin1991}. In conventional Finite Element Methods (FEM), node-based scalar approximations that interpolate Cartesian vector components ($E_x, E_y$) independently suffer from three well-known pathologies \cite{jin2014}: they impose artificial $C^0$ continuity on normal field components across material interfaces, fail to resolve field singularities at sharp conducting corners and edges, and violate the discrete de Rham sequence, contaminating the discrete spectrum with non-physical irrotational spurious modes. N\'ed\'elec's edge finite elements (EFEM) overcome these flaws by associating degrees of freedom with tangential field circulations along element edges, natively conforming in $H(\mathbf{curl})$; however, they remain heavily constrained by the generation and adaptation of high-quality conforming meshes, which becomes a major bottleneck for curved inclusions and reentrant geometries.

To eliminate mesh-connectivity constraints, meshless methods construct field approximations directly over scattered point clouds \cite{belytschko1994}. Nevertheless, early vector adaptations of scalar meshless schemes (such as Element-Free Galerkin and Local Petrov-Galerkin methods) inherited the same drawbacks as nodal FEM---spectral pollution and normal-jump locking---requiring auxiliary divergence-free projections or Lagrange multipliers. Seeking native $H(\mathbf{curl})$ regularity without a mesh, Lima and Mesquita \cite{lima2017} introduced the Edge Meshless Method (EMM), formulating vector shape functions from field circulations along directed edge segments. By taking the asymptotic limit as edge lengths shrink to zero ($L_k \to 0$), Ortiz et al. \cite{ortiz2021,ortiz2023} established the \textbf{Vector Nodal Meshless Method (VNMM)}, in which each edge degenerates into a vector node $(\mathbf{x}_i, \mathbf{t}_i)$ carrying a single scalar directional projection degree of freedom, $e_i = \mathbf{E}(\mathbf{x}_i)\cdot\mathbf{t}_i$.

By employing a complete linear polynomial vector basis ($\mathcal{P}^1$) with $k$-d tree support selection, VNMM avoids the curl-convergence stagnation typical of incomplete solenoidal bases ($\mathcal{L}^1$) and preserves exact first-order polynomial completeness on arbitrary point distributions. Furthermore, because dynamic local stencils on unconstrained point clouds break global trace conformity, incorporating variational div-curl regularization ($s_{\mathrm{div}}\mathbf{K}_{\mathrm{div}}$) shifts spurious gradient modes outside the operational spectrum without perturbing physical resonances.

Despite these advances in homogeneous domains, extending VNMM to \textbf{heterogeneous, discontinuous, and anisotropic media} introduces a critical challenge: sharing support stencils across distinct dielectrics smears the physical normal discontinuity ($\varepsilon_1 E_{1n} = \varepsilon_2 E_{2n}$) and induces Gibbs oscillations. This work resolves that limitation through a \textbf{strict material support partition rule} combined with \textbf{tangential interface directors} ($\mathbf{t}_j = \boldsymbol{\tau}$), coupling adjacent media via a single tangential degree of freedom while leaving the normal component free to jump sharply. Moreover, we demonstrate that VNMM naturally thrives on \textbf{unstructured irregular point clouds}: breaking Cartesian grid symmetry via moderate stochastic jitter (10\%--25\%) mitigates directional dispersion (\textit{grid-locking}) and reduces mean eigenvalue errors by over 50\%, while accurately resolving complex Hermitian gyrotropic tensors, reentrant corner singularities, curved PEC inclusions, and multi-material intersections without spurious modes.

The remainder of this paper is organized as follows. Section~2 presents the 2D VNMM mathematical formulation on point clouds, the scale-invariant automated support selection algorithm, the strict multimaterial support partition conditions, and the regularized Galerkin weak form. Section~3 investigates polynomial completeness, stochastic jitter sensitivity, and sharp interface jump resolution. Section~4 validates dielectric-loaded cavities across wide permittivity contrasts. Section~5 evaluates benchmarks in anisotropic waveguides, corner singularities, curved PEC inclusions, and a four-material junction, followed by concluding remarks in Section~6.

\section{Mathematical Formulation of the Multimaterial 2D VNMM on Point Clouds}

This section establishes the mathematical formulation of the two-dimensional Vector Nodal Meshless Method (2D VNMM) for heterogeneous and anisotropic media on unstructured point clouds. First, we define the directional projection degrees of freedom, the complete linear polynomial vector basis ($\mathcal{P}^1$), and the local collocation system. Next, we describe the scale-invariant automated support selection algorithm coupled with strict multimaterial partition conditions, followed by the divergence-regularized weak Galerkin formulation.

\subsection{Degrees of Freedom, Nodal Vector Projections, and PEC Boundaries}
Consider a two-dimensional domain $\Omega \subset \mathbb{R}^2$ discretized by a cloud of $N$ arbitrary nodes $\{\mathbf{x}_i\}_{i=1}^N$ distributed throughout the interior and boundaries $\partial\Omega$. Unlike scalar meshless methods that assign separate degrees of freedom to Cartesian components ($E_x, E_y$), VNMM associates each node $\mathbf{x}_i = (x_i, y_i)^T$ with a single unit director vector:
\begin{equation}
\mathbf{t}_i = \begin{bmatrix} \cos\theta_i \\ \sin\theta_i \end{bmatrix}
\end{equation}

At any evaluation point $\mathbf{x} = (x, y) \in \Omega$, the approximated electric field $\mathbf{E}^h(\mathbf{x})$ is expressed as a linear combination of local vector nodal shape functions $\mathbf{N}_i(\mathbf{x})$ associated with the $M_{\mathrm{supp}}$ nodes in the local support stencil:
\begin{equation}
\mathbf{E}^h(x, y) = \sum_{i=1}^{M_{\mathrm{supp}}} \mathbf{N}_i(x, y)\, e_i
\end{equation}
where $e_i$ denotes the scalar nodal degree of freedom representing the pointwise directional projection of the electric field along the nodal director:
\begin{equation}
e_i = \mathbf{E}(\mathbf{x}_i) \cdot \mathbf{t}_i \quad [\text{V/m}]
\end{equation}

To ensure local $H(\mathbf{curl})$ consistency and direct enforcement of essential boundary conditions, the vector shape functions are required to satisfy the vector Kronecker delta projection property:
\begin{equation}
\mathbf{N}_i(\mathbf{x}_k) \cdot \mathbf{t}_k = \delta_{ik} = \begin{cases} 1, & \text{if } k = i \\ 0, & \text{if } k \neq i \end{cases}
\end{equation}

Along Perfect Electric Conductor (PEC) boundaries, where $\hat{n} \times \mathbf{E} = \mathbf{0}$, director vectors are aligned with the boundary tangent ($\mathbf{t}_k = \boldsymbol{\tau}$), allowing homogeneous Dirichlet conditions to be enforced directly via nodal collocation ($e_k = 0, \; \forall \mathbf{x}_k \in \partial\Omega_{\mathrm{PEC}}$) without Lagrange multipliers. In the interior of each subdomain, an alternating orthogonal checkerboard pattern ($\pm 45^\circ$, i.e., $\theta_i \in \{\pi/4, 3\pi/4\}$) is assigned to ensure local angular diversity. A detailed functional analysis of this directional projection scheme and the distinction between pointwise nodal collocation and continuous boundary trace on unconstrained point clouds is presented in \cite{mesquita_eabe}.

\subsection{Complete Linear Polynomial Vector Basis (\texorpdfstring{$\mathcal{P}^1$}{P1})}
Early VNMM formulations employed the incomplete 3-node solenoidal space $\mathcal{L}^1$ inherited from lowest-order N\'ed\'elec edge elements. However, when combined with dynamic support selection on arbitrary point clouds, the $\mathcal{L}^1$ basis artificially forces longitudinal field derivatives ($\partial E_x/\partial x$, $\partial E_y/\partial y$) to vanish, triggering derivative modal aliasing and $\mathcal{O}(1)$ curl-convergence stagnation.

To overcome this limitation, we adopt the complete linear polynomial vector space $\mathcal{P}^1 = \mathcal{P}_1 \times \mathcal{P}_1$, comprising 6 independent monomials ($M_{\mathrm{supp}} = 6$ support nodes) in local translated coordinates ($\Delta x = x - x_P$, $\Delta y = y - y_P$) centered at the evaluation point $P(x_P, y_P)$:
\begin{equation}
\mathcal{P}^1 = \operatorname{span}\left\{
\begin{bmatrix} 1 \\ 0 \end{bmatrix}, \;
\begin{bmatrix} 0 \\ 1 \end{bmatrix}, \;
\begin{bmatrix} \Delta x \\ 0 \end{bmatrix}, \;
\begin{bmatrix} \Delta y \\ 0 \end{bmatrix}, \;
\begin{bmatrix} 0 \\ \Delta x \end{bmatrix}, \;
\begin{bmatrix} 0 \\ \Delta y \end{bmatrix}
\right\}
\end{equation}
The local vector shape function associated with support node $i$ is written as:
\begin{equation}
\mathbf{N}_i(\Delta x, \Delta y) = \begin{bmatrix} \beta_{1i} + \beta_{3i}\Delta x + \beta_{4i}\Delta y \\ \beta_{2i} + \beta_{5i}\Delta x + \beta_{6i}\Delta y \end{bmatrix}
\end{equation}

With 6 independent degrees of freedom, the $\mathcal{P}^1$ space resolves all four components of the local Jacobian tensor $\mathbf{J}^h = \nabla \mathbf{E}^h$ simultaneously, capturing both rotational and longitudinal flux variations. A complete proof that this transition eliminates derivative aliasing and restores optimal convergence of the field and its curl is provided in \cite{mesquita_eabe}.

\subsection{Local Collocation System and Differential Operators}
For a qualified set of 6 neighboring support nodes $\mathcal{S}(P) = \{n_1, \dots, n_6\}$, enforcing the vector Kronecker delta condition ($\mathbf{N}_i(\mathbf{x}_k) \cdot \mathbf{t}_k = \delta_{ik}$) yields the local algebraic placement system:
\begin{equation}
\mathbf{A}\boldsymbol{\beta} = \mathbf{I}_6 \implies \boldsymbol{\beta} = \mathbf{A}^{-1}
\end{equation}
where the $k$-th row of the placement matrix $\mathbf{A} \in \mathbb{R}^{6 \times 6}$ is given by:
\begin{equation}
\mathbf{A}_{k,:} = \begin{bmatrix} t_{kx} & t_{ky} & \Delta x_k t_{kx} & \Delta y_k t_{kx} & \Delta x_k t_{ky} & \Delta y_k t_{ky} \end{bmatrix}, \quad k = 1, \dots, 6
\end{equation}
with $\Delta x_k = x_k - x_P$ and $\Delta y_k = y_k - y_P$.

Evaluating directly at point $P$ ($\Delta x = \Delta y = 0$), the vector shape functions and their spatial differential operators are obtained analytically from the rows of $\boldsymbol{\beta}$ (0-indexed):
\begin{align}
\boldsymbol{\Phi}(P) &= \boldsymbol{\beta}_{0:2, :} \in \mathbb{R}^{2 \times 6} \\
(\nabla \times \boldsymbol{\Phi})_z(P) &= \boldsymbol{\beta}_{4, :} - \boldsymbol{\beta}_{3, :} \in \mathbb{R}^{1 \times 6} \\
\nabla \cdot \boldsymbol{\Phi}(P) &= \boldsymbol{\beta}_{2, :} + \boldsymbol{\beta}_{5, :} \in \mathbb{R}^{1 \times 6}
\end{align}
The algebraic characterization of singularities in $\mathbf{A}$ via the Null-Space Orthogonality Principle and its dimensionless factorization $\mathbf{A} = \widehat{\mathbf{A}}\mathbf{D}_h$ to preserve matrix conditioning under spatial refinement ($h \to 0$) are rigorously established in \cite{mesquita_eabe}.

\subsection{Automated Support Selection and Strict Material Partition Rule}
Selecting a valid support stencil $\mathcal{S}(P)$ on unstructured point clouds requires ensuring both algebraic well-conditioning and strict respect for material discontinuities.

\paragraph{Scale-Invariant Support Selection Algorithm.}
On arbitrary point clouds, candidate sextets are selected using the 4-phase automated support selection algorithm introduced in \cite{mesquita_eabe}:
\begin{enumerate}
  \item \textbf{Spatial Neighborhood Query via $k$-d Tree:} An orthogonal $k$-d tree ($\mathcal{O}(N \log N)$) retrieves the $K_{\mathrm{initial}} \ge 12$ nearest neighbors of $P$, defining the local characteristic spacing $h_{\mathrm{local}}(\mathbf{x}_P) = \frac{1}{6}\sum_{k=1}^6 \|\mathbf{x}_{(k)} - \mathbf{x}_P\|$.
  \item \textbf{Proximity-Ordered Candidate Generation:} Anchoring the nearest node to $P$, candidate 6-node combinations are generated in order of increasing mean distance.
  \item \textbf{Dimensionless Matrix Assembly:} Local coordinates are normalized ($\hat{x}_k = \Delta x_k/h_{\mathrm{local}}, \; \hat{y}_k = \Delta y_k/h_{\mathrm{local}}$) to assemble the dimensionless matrix $\widehat{\mathbf{A}}$, eliminating quartic determinant decay ($\det(\mathbf{A}) \sim \mathcal{O}(h^4)$) and artificial condition-number inflation ($\kappa_2(\mathbf{A}) \sim \mathcal{O}(1/h)$).
  \item \textbf{Determinant Qualification with Early Stopping:} The first candidate sextet satisfying $|\det(\widehat{\mathbf{A}})| \ge \widehat{\mathrm{Tol}} > 0$ is immediately accepted, screening out singular or ill-conditioned stencils without support dilation.
\end{enumerate}

\paragraph{Conditions for Multimaterial Domains (Strict Support Partition Rule).}
For heterogeneous domains composed of distinct materials ($\Omega = \bigcup_{m=1}^M \Omega_m$) separated by interfaces $\Gamma_{\mathrm{int}}$, the $k$-d tree neighborhood search must strictly enforce two additional physical conditions:
\begin{enumerate}
  \item \textbf{Subdomain Confinement:} When evaluating shape functions at an integration point $P \in \Omega_m$, candidate support nodes are strictly restricted to nodes belonging to medium $\Omega_m$ and nodes lying directly on the adjacent interface $\Gamma_{\mathrm{int}}$ (implemented via material-partitioned $k$-d trees). Nodes from other media are strictly excluded from the search.
  \item \textbf{Tangential Interface Degrees of Freedom:} Nodes located along physical interfaces $\Gamma_{\mathrm{int}}$ have their director vectors aligned exclusively with the interface tangent ($\mathbf{t}_j = \boldsymbol{\tau}$). Because the tangential electric field is physically continuous ($E_{1t} = E_{2t}$), a single scalar degree of freedom on the interface couples both media exactly. Meanwhile, support segregation allows the normal component $E_n$ to undergo the natural physical step jump governed by the permittivity ratio ($\varepsilon_1 E_{1n} = \varepsilon_2 E_{2n}$) without artificial smearing or Gibbs oscillations.
\end{enumerate}

\subsection{Weak Galerkin Formulation with Divergence Regularization}
Because dynamic stencil selection on irregular point clouds yields a globally non-conforming approximation ($V_h \not\subset H(\mathbf{curl}; \Omega)$), variational div-curl regularization \cite{mesquita_eabe} is incorporated into the generalized Maxwell eigenvalue weak form for arbitrary anisotropic media characterized by permittivity tensor $\overline{\overline{\varepsilon}}_r$ and permeability $\mu_r$:
\begin{equation}
\int_{\Omega} \frac{1}{\mu_r} (\nabla \times \mathbf{E})_z (\nabla \times \mathbf{v})_z^* \, d\Omega + s_{\mathrm{div}} \int_{\Omega} \frac{1}{\mu_r} (\nabla \cdot \mathbf{E}) (\nabla \cdot \mathbf{v})^* \, d\Omega = k_0^2 \int_{\Omega} \mathbf{v}^* \cdot (\overline{\overline{\varepsilon}}_r \mathbf{E}) \, d\Omega
\end{equation}
where $s_{\mathrm{div}} \ge 0$ is the divergence regularization penalty parameter that shifts irrotational spurious modes out of the physical spectrum, and $(\cdot)^*$ denotes the conjugate transpose for complex Hermitian tensors.

Discretization leads to the generalized eigenvalue problem:
\begin{equation}
(\mathbf{K}_{\mathrm{curl}} + s_{\mathrm{div}} \mathbf{K}_{\mathrm{div}}) \mathbf{c} = \lambda \mathbf{M} \mathbf{c}
\end{equation}
where $\lambda = k_0^2$. Homogeneous Dirichlet boundary conditions on PEC walls ($E_t = 0$) are applied directly by eliminating the corresponding boundary rows and columns ($e_k = 0, \; \forall k \in \partial\Omega_{\mathrm{PEC}}$). Numerical integration is performed on a uniform background Cartesian cell grid using Gauss quadrature ($2 \times 2$, or $3 \times 3$ near singular vertices), serving strictly as a metric support without requiring nodal connectivity.

\section{Consistency and Robustness Study on Irregular Point Clouds}

This section systematically evaluates the mathematical consistency, spatial convergence, and numerical stability of the multimaterial 2D VNMM on irregular point clouds, covering polynomial reproduction tests (Patch Test and MMS), statistical sensitivity to stochastic nodal jitter up to 50\%, and sharp resolution of normal field jumps across dielectric interfaces.

\subsection{Polynomial Completeness and Spatial Convergence}
To verify the consistency of VNMM shape functions on multimaterial point clouds, the Galerkin Patch Test and Method of Manufactured Solutions (MMS) were conducted:
\begin{itemize}
  \item \textbf{Affine Patch Test Field:}
  \begin{equation}
  \mathbf{E}_{\mathrm{patch}}(x, y) = \begin{bmatrix} 1.5 + 0.8\,x - 0.3\,y \\ -0.7 + 0.4\,x + 1.2\,y \end{bmatrix}, \quad (\nabla \times \mathbf{E}_{\mathrm{patch}})_z = 0.7\text{ V/m}^2
  \end{equation}
  \item \textbf{Non-Polynomial MMS Field:}
  \begin{equation}
  \mathbf{E}_{\mathrm{MMS}}(x, y) = \begin{bmatrix} \sin\!\left(\frac{\pi x}{L_x}\right)\cos\!\left(\frac{\pi y}{L_y}\right) \\ -\cos\!\left(\frac{\pi x}{L_x}\right)\sin\!\left(\frac{\pi y}{L_y}\right) \end{bmatrix}, \quad (\nabla \times \mathbf{E}_{\mathrm{MMS}})_z = \left(\frac{\pi}{L_x} + \frac{\pi}{L_y}\right)\sin\!\left(\frac{\pi x}{L_x}\right)\sin\!\left(\frac{\pi y}{L_y}\right)
  \end{equation}
\end{itemize}

The quantitative interpolation errors and asymptotic spatial convergence rates are summarized in Table~\ref{tab:interp}, while Figure~\ref{fig:interp} illustrates the spatial error distribution across the bimaterial domain, machine-precision Patch Test reproduction, continuous reconstructed vector fields, and the geometric verification of support stencils strictly partitioned across material regions.

\begin{table}[H]
\centering
\caption{VNMM 2D interpolation performance on a bimaterial domain using point clouds.}
\label{tab:interp}
\resizebox{\textwidth}{!}{
\begin{tabular}{llccc}
\toprule
Test & Vector Field & Relative $L_2$ Error ($\mathbf{E}$) & Relative Error $\mathrm{curl}(\mathbf{E})$ & Status \\
\midrule
\textbf{Galerkin Patch Test} & Complete Affine Linear & $\mathbf{2.85 \times 10^{-16}}$ & $\mathbf{1.00 \times 10^{-14}}$ & \textbf{Machine Precision} \\
\textbf{MMS (Base $N=9$)} & Smooth Trigonometric & $5.742\%$ & $19.461\%$ & Consistent \\
\textbf{MMS (Refined $N=17$)} & Smooth Trigonometric & $\mathbf{1.147\%}$ & $\mathbf{9.390\%}$ & \textbf{$\mathcal{O}(h^2)$ Convergence} \\
\textbf{Asymptotic Rate} & --- & $\mathbf{5.01\times}$ ($\ge 4\times$) & $\mathbf{2.07\times}$ ($\approx 2\times$) & \textbf{Optimal Theoretical} \\
\bottomrule
\end{tabular}}
\end{table}

\begin{figure}[H]
\centering
\includegraphics[width=0.95\textwidth]{VNMM_Materials_interpolation_results.png}
\caption{Interpolation panel across a bimaterial domain: machine-precision patch test reproduction (top-left), MMS error distribution (top-right), reconstructed vector field (bottom-left), and geometric verification of support stencils strictly partitioned across material regions (bottom-right).}
\label{fig:interp}
\end{figure}

The results in Table~\ref{tab:interp} and Figure~\ref{fig:interp} confirm that the exact first-order completeness and polynomial reproducing properties ($\mathcal{P}^1$) are fully inherited by the multimaterial formulation, reaching the machine precision floor ($10^{-16}$). This rigorously proves that the strict domain-based support partitioning and the assignment of tangential interface directors introduce no spurious distortions and do not compromise the fundamental consistency and reproduction properties established in homogeneous media. Furthermore, the MMS test confirms that the presence of heterogeneous interfaces preserves the optimal asymptotic spatial convergence rates: quadratic for the electric field ($\mathcal{O}(h^2)$) and linear for its curl ($\mathcal{O}(h)$), mathematically certifying the robustness and consistency of VNMM in piecewise continuous media.

\subsection{Statistical Sensitivity Analysis under Nodal Jitter (0\% to 50\%)}
To demonstrate the genuine meshless nature of VNMM and its robustness to spatial irregularity, a random stochastic perturbation was applied to all interior nodes:
\begin{equation}
\mathbf{x}_i \leftarrow \mathbf{x}_i + \eta \, h_{\min} \begin{bmatrix} \xi_x \\ \xi_y \end{bmatrix}, \quad \xi_x, \xi_y \sim \mathcal{U}(-0.5, 0.5)
\end{equation}
where $\eta \in \{0.00, 0.10, 0.25, 0.40, 0.50\}$ denotes the jitter fraction, evaluated across multiple random seeds ($N_{\mathrm{seeds}} = 3$) on the dielectric-loaded cavity ($[0, \pi]^2$, interface at $x = \pi/2$, $\varepsilon_{r1}=1.0$, $\varepsilon_{r2}=4.0$).

The global statistical performance across jitter levels is summarized in Table~\ref{tab:jitter}, while Table~\ref{tab:jitter_modes} details the cutoff wavenumbers and relative errors of the first 6 resonant modes against analytical roots. In addition, Figure~\ref{fig:jitter} illustrates the geometry of the perturbed point cloud under 25\% displacement, the relative error curves in $k_c$ with statistical dispersion bars ($\mu \pm \sigma$), and the continuous monitoring of residual divergence energy ratio ($R_{\mathrm{div}} < 3.3\%$) verifying spectral purity.

\begin{table}[H]
\centering
\caption{Overall statistical performance of 2D VNMM under stochastic nodal jitter up to 50\%.}
\label{tab:jitter}
\resizebox{\textwidth}{!}{
\begin{tabular}{ccccccc}
\toprule
Jitter $\eta$ & Patch Test Error ($L_2^{\mathrm{rel}}$) & MMS Error ($L_2^{\mathrm{rel}}$) & Mode 1 $k_c$ ($\mu \pm \sigma$) & Mode 1 Error (\%) & Mean Error (6 Modes) (\%) & Residual $R_{\mathrm{div}}$ (\%) \\
\midrule
\textbf{0\% (Regular)} & $2.35 \times 10^{-16}$ & $3.22\%$ & $0.5825 \pm 0.000$ & $4.23\%$ & $2.46 \pm 0.00\%$ & $1.28\%$ \\
\textbf{10\% (Mild)} & $3.53 \times 10^{-16}$ & $3.66\%$ & $0.6037 \pm 0.002$ & $\mathbf{0.75 \pm 0.35\%}$ & $\mathbf{1.13 \pm 0.21\%}$ & $1.11\%$ \\
\textbf{25\% (Moderate)} & $4.85 \times 10^{-16}$ & $6.92\%$ & $0.6226 \pm 0.005$ & $\mathbf{2.37 \pm 0.85\%}$ & $\mathbf{1.22 \pm 0.17\%}$ & $1.71\%$ \\
\textbf{40\% (Severe)} & $3.22 \times 10^{-16}$ & $5.67\%$ & $0.6365 \pm 0.002$ & $4.65 \pm 0.39\%$ & $3.32 \pm 0.12\%$ & $2.71\%$ \\
\textbf{50\% (Extreme)} & $5.24 \times 10^{-16}$ & $7.97\%$ & $0.6463 \pm 0.001$ & $6.27 \pm 0.10\%$ & $5.27 \pm 0.26\%$ & $3.26\%$ \\
\bottomrule
\end{tabular}}
\end{table}

\begin{table}[H]
\centering
\caption{Cutoff wavenumbers $k_c$ (rad/m) and relative errors of the first 6 resonant modes computed by 2D VNMM across stochastic jitter levels ($\eta \in [0\%, 50\%]$).}
\label{tab:jitter_modes}
\resizebox{\textwidth}{!}{
\begin{tabular}{ccccccc}
\toprule
Mode & Exact Analytical & $\eta = 0\%$ (Regular) & $\eta = 10\%$ (Mild) & $\eta = 25\%$ (Moderate) & $\eta = 40\%$ (Severe) & $\eta = 50\%$ (Extreme) \\
\midrule
\textbf{1} & $0.6082$ & $0.5824$ ($4.23\%$) & $0.6036$ ($0.75\%$) & $0.5937$ ($2.37\%$) & $0.5799$ ($4.65\%$) & $0.5701$ ($6.27\%$) \\
\textbf{2} & $0.6421$ & $0.6018$ ($6.28\%$) & $0.6425$ ($0.06\%$) & $0.6342$ ($1.24\%$) & $0.6194$ ($3.54\%$) & $0.6088$ ($5.19\%$) \\
\textbf{3} & $1.0385$ & $1.0370$ ($0.14\%$) & $1.0399$ ($0.14\%$) & $1.0269$ ($1.12\%$) & $1.0097$ ($2.78\%$) & $0.9923$ ($4.45\%$) \\
\textbf{4} & $1.0992$ & $1.1119$ ($1.16\%$) & $1.1180$ ($1.71\%$) & $1.0978$ ($0.13\%$) & $1.0724$ ($2.44\%$) & $1.0509$ ($4.40\%$) \\
\textbf{5} & $1.3918$ & $1.4092$ ($1.25\%$) & $1.3928$ ($0.07\%$) & $1.3829$ ($0.64\%$) & $1.3452$ ($3.35\%$) & $1.3065$ ($6.13\%$) \\
\textbf{6} & $1.5328$ & $1.5590$ ($1.71\%$) & $1.5698$ ($2.42\%$) & $1.5481$ ($1.00\%$) & $1.4843$ ($3.17\%$) & $1.4537$ ($5.16\%$) \\
\bottomrule
\end{tabular}}
\end{table}

\begin{figure}[H]
\centering
\includegraphics[width=\textwidth]{VNMM_Jitter_Detailed_Study.png}
\caption{Statistical study of sensitivity to irregular point clouds (0\% to 50\% jitter): (a) Geometric configuration of the nodal cloud under 25\% jitter; (b) Cutoff wavenumber relative error with statistical dispersion bars ($\mu \pm \sigma$); (c) Continuous monitoring of residual divergence energy ratio ($R_{\mathrm{div}} < 3.3\%$) and MMS interpolation error.}
\label{fig:jitter}
\end{figure}

These quantitative and visual results reveal three key physical and numerical insights:
\begin{enumerate}
  \item \textbf{Exact Invariance of Order-1 Consistency:} The vector Patch Test error remained strictly at $10^{-16}$ across \textbf{all perturbation levels}, from fully regular (0\%) to highly disordered (50\%), confirming that reproducing polynomial fields is independent of Cartesian point alignment.
  \item \textbf{Mitigation of Grid-Locking Effects:} On rigid regular points (0\%), preferred grid directions cause anisotropic dispersion, yielding a $2.46\%$ mean spectral error. Moderate irregular jitter ($\eta = 10\%$ and $25\%$) reduced the mean eigenvalue error to $\mathbf{1.13\%}$ and $\mathbf{1.22\%}$, respectively --- a \textbf{performance gain exceeding 50\%} due to breaking Cartesian grid symmetry.
  \item \textbf{Spectral Purity under High Deformation:} Even under 50\% severe jitter, the mean residual divergence energy remained bounded at $R_{\mathrm{div}} = 3.26\%$ (well below the 5.0\% threshold), ensuring total suppression of spurious modes.
\end{enumerate}

\subsection{Sharp Capture of Normal Jump Discontinuities (\texorpdfstring{$[E_n] \neq 0$}{[En] != 0})}
To assess the local field behavior across material interfaces, a 1D transverse cut at $y = 0.5\text{ m}$ was evaluated for a field exhibiting an exact $4:1$ physical normal jump ($E_{1x} = 1.500\text{ V/m} \to E_{2x} = 0.375\text{ V/m}$), while the tangential component remains strictly continuous ($E_{1y} = E_{2y} = 0.500\text{ V/m}$).

The spectral stability of the first 6 resonant modes in the bimaterial cavity under nodal perturbation ($\eta \in \{0\%, 10\%, 20\%\}$) is summarized in Table~\ref{tab:jump_stability}, while Figure~\ref{fig:jump} illustrates the 1D transverse profiles capturing the sharp 4:1 normal jump without Gibbs oscillations, strict $C^0$ tangential continuity across the interface, and the fundamental mode tracking the exact analytical dispersion curve across permittivity contrasts ($\varepsilon_{r2} \in [2.25, 16.0]$).

\begin{table}[H]
\centering
\caption{Spectral performance of the first 6 resonant modes in the bimaterial cavity under nodal perturbation ($\eta \in \{0\%, 10\%, 20\%\}$).}
\label{tab:jump_stability}
\resizebox{\textwidth}{!}{
\begin{tabular}{cccccc}
\toprule
Mode & Classification & Exact Analytical $k_c$ (rad/m) & $\eta = 0\%$ (Regular) & $\eta = 10\%$ (Mild) & $\eta = 20\%$ (Moderate) \\
\midrule
\textbf{1} & $\mathrm{TE}(0,1)$ & $0.6082$ & $0.6018$ ($1.04\%$) & $0.6068$ ($0.23\%$) & $0.6022$ ($0.98\%$) \\
\textbf{2} & $\mathrm{TE}(1,1)$ & $0.6421$ & $0.5824$ ($9.30\%$) & $0.6474$ ($0.82\%$) & $0.6394$ ($0.43\%$) \\
\textbf{3} & $\mathrm{TE}(0,2)$ & $1.0385$ & $1.0370$ ($0.14\%$) & $1.0426$ ($0.40\%$) & $1.0285$ ($0.96\%$) \\
\textbf{4} & $\mathrm{TE}(1,2)$ & $1.0992$ & $1.1119$ ($1.16\%$) & $1.1226$ ($2.12\%$) & $1.1013$ ($0.19\%$) \\
\textbf{5} & $\mathrm{TE}(2,1)$ & $1.3918$ & $1.4092$ ($1.25\%$) & $1.4123$ ($1.47\%$) & $1.4036$ ($0.84\%$) \\
\textbf{6} & $\mathrm{TE}(0,3)$ & $1.5328$ & $1.5590$ ($1.71\%$) & $1.5790$ ($3.02\%$) & $1.5694$ ($2.39\%$) \\
\textbf{Mean Error} & --- & --- & $\mathbf{2.43\%}$ & $\mathbf{1.34\%}$ & $\mathbf{0.97\%}$ \\
\bottomrule
\end{tabular}}
\end{table}

\begin{figure}[H]
\centering
\includegraphics[width=\textwidth]{VNMM_Interface_Jump_and_Robustness.png}
\caption{VNMM performance across dielectric interfaces and permittivity contrast: (a) 1D transverse profile of $E_x(x)$ demonstrating a sharp 4:1 normal jump without Gibbs oscillations; (b) 1D transverse profile of $E_y(x)$ verifying strict $C^0$ continuity; (c) Fundamental cutoff wavenumber vs dielectric contrast ($\varepsilon_{r2} \in [2.25, 16.0]$).}
\label{fig:jump}
\end{figure}

The results in Table~\ref{tab:jump_stability} and Figure~\ref{fig:jump} lead to three important conclusions regarding interface resolution:
\begin{enumerate}
  \item \textbf{Sharp Normal Jump without Smearing:} The 1D transverse profile $E_x(x)$ captures the 4:1 normal discontinuity with high fidelity (maximum error $< 0.033\text{ V/m}$) exactly at $x = 1.0\text{ m}$, completely free of numerical smearing or Gibbs oscillations;
  \item \textbf{Conformity in $H(\mathbf{curl})$ and Tangential Continuity:} Strict physical $C^0$ continuity is preserved across the interface in $E_y(x)$ (maximum error $< 0.014\text{ V/m}$), confirming that shared tangential degrees of freedom enforce proper interface boundary conditions;
  \item \textbf{Spectral Robustness under Nodal Irregularity:} Nodal jitter reduces the average error across the first 6 modes from $2.43\%$ to $0.97\%$, demonstrating that the method accurately models multimaterial interfaces even on irregular point clouds.
\end{enumerate}

\subsection{Synthesis and Summary of Consistency and Robustness Analysis}
The numerical benchmarks presented in this section establish three foundational properties of 2D VNMM in multimaterial domains:
\begin{enumerate}
  \item \textbf{Invariant Completeness and Polynomial Reproduction:} The directional-projection formulation preserves exact first-order completeness ($\mathcal{P}^1$) down to machine precision ($10^{-16}$) across arbitrary point distributions, exhibiting complete robustness against extreme stochastic nodal disorder (up to 50\% jitter);
  \item \textbf{Mitigation of Anisotropic Numerical Dispersion:} Moderately irregular point clouds (10\% to 25\% jitter) mitigate Cartesian grid-locking, cutting average eigenvalue errors by over 50\% relative to rigid regular grids while strictly maintaining spectral purity ($R_{\mathrm{div}} < 3.3\%$);
  \item \textbf{Exact Resolution of Interface Jump Conditions:} Domain-segregated local stencils coupled with tangential interface directors decouple discontinuous normal components $[E_n] \neq 0$ while enforcing tangential continuity $[E_t] = 0$ without numerical smearing, demonstrating the viability of vector nodal meshless methods for heterogeneous electromagnetic systems.
\end{enumerate}

\section{Eigenvalue Solutions in Cavities and Waveguides}

\subsection{Dielectric-Loaded Cavity: Transcendental Analytical Comparison}
The rectangular cavity $[0, \pi]^2$ with an interface at $x = \pi/2$ ($\varepsilon_{r1} = 1.0$, $\varepsilon_{r2} = 4.0$) was benchmarked against the exact roots of transcendental dispersion equations:
\begin{align}
k_{x1} \cot(k_{x1} d) + k_{x2} \cot(k_{x2} (a - d)) &= 0 \quad (m = 0) \\
\frac{\beta_1}{\varepsilon_{r1}} \tan(\beta_1 d) + \frac{\beta_2}{\varepsilon_{r2}} \tan(\beta_2 (a - d)) &= 0 \quad (m \ge 1)
\end{align}

The quantitative comparison between the 2D VNMM eigenvalues and exact analytical solutions for the first 8 resonant $\mathrm{TE}_z$ modes is reported in Table~\ref{tab:cavity}, while Figure~\ref{fig:cavity_modes} displays the normalized spatial electric field magnitude ($\|\mathbf{E}\|/\|\mathbf{E}\|_{\max}$) and white vector arrows for the first four modes. In particular, the vector arrows in Figure~\ref{fig:cavity_modes} highlight both the energy confinement inside the dielectric slab ($\Omega_2$) and the sharp discontinuity of the normal field component ($E_x$) across $x = \pi/2$---clearly visible in Modes 2, 3, and 4 due to the $4:1$ permittivity contrast---while the tangential component ($E_y$) remains continuous.

\begin{table}[H]
\centering
\caption{Comparison between 2D VNMM eigenvalues and exact analytical solutions for the first 8 resonant $\mathrm{TE}_z$ modes.}
\label{tab:cavity}
\resizebox{\textwidth}{!}{
\begin{tabular}{cccccccc}
\toprule
Mode & Classification & Analytical $\lambda$ & VNMM 2D $\lambda$ & Analytical $k_c$ (rad/m) & VNMM $k_c$ (rad/m) & Rel. $k_c$ Error (\%) & $R_{\mathrm{div}}$ (\%) \\
\midrule
\textbf{1} & $\mathrm{TE}(0,1)$ & $0.369888$ & $0.364022$ & $0.608184$ & $0.603342$ & $\mathbf{0.80\%}$ & $0.373\%$ \\
\textbf{2} & $\mathrm{TE}(1,1)$ & $0.412351$ & $0.416827$ & $0.642146$ & $0.645621$ & $\mathbf{0.54\%}$ & $0.083\%$ \\
\textbf{3} & $\mathrm{TE}(0,2)$ & $1.078512$ & $1.042186$ & $1.038514$ & $1.020875$ & $\mathbf{1.70\%}$ & $0.808\%$ \\
\textbf{4} & $\mathrm{TE}(1,2)$ & $1.208151$ & $1.155490$ & $1.099159$ & $1.074937$ & $\mathbf{2.20\%}$ & $0.852\%$ \\
\textbf{5} & $\mathrm{TE}(2,1)$ & $1.937222$ & $1.980649$ & $1.391841$ & $1.407355$ & $\mathbf{1.11\%}$ & $0.169\%$ \\
\textbf{6} & $\mathrm{TE}(0,3)$ & $2.349479$ & $2.180424$ & $1.532801$ & $1.476626$ & $\mathbf{3.66\%}$ & $0.380\%$ \\
\textbf{7} & $\mathrm{TE}(1,3)$ & $2.583093$ & $2.500588$ & $1.607200$ & $1.581325$ & $\mathbf{1.61\%}$ & $0.574\%$ \\
\textbf{8} & $\mathrm{TE}(2,2)$ & $2.747205$ & $2.877572$ & $1.657469$ & $1.696341$ & $\mathbf{2.35\%}$ & $1.809\%$ \\
\bottomrule
\end{tabular}}
\end{table}

\begin{figure}[H]
\centering
\includegraphics[width=0.92\textwidth]{VNMM_Dielectric_Cavity_Modes.png}
\caption{Normalized field magnitude $\|\mathbf{E}\|/\|\mathbf{E}\|_{\max}$ and electric field vectors (white arrows) for the first four resonant modes in the dielectric-loaded cavity, highlighting energy concentration inside the higher-permittivity slab ($\varepsilon_{r2}=4.0$) and the normal field discontinuity at the interface in Modes 2, 3, and 4.}
\label{fig:cavity_modes}
\end{figure}

The results in Table~\ref{tab:cavity} and Figure~\ref{fig:cavity_modes} confirm that 2D VNMM achieved an average error of only $1.75\%$ across the entire spectrum (maximum $3.66\%$), with divergence energy ratio $R_{\mathrm{div}} < 2.0\%$ for all physical modes, proving complete suppression of spurious spectral pollution and faithful reproduction of interface jump physics.

\subsection{Permittivity Contrast Parametric Sweep}
The spectral response was evaluated across practical microwave substrate permittivities ($\Omega_1$ air with $\varepsilon_{r1} = 1.0$, $\Omega_2$ dielectric): Teflon ($\varepsilon_{r2} = 2.25$), Standard Dielectric ($\varepsilon_{r2} = 4.00$), Alumina ($\varepsilon_{r2} = 9.80$), and Advanced Ceramic ($\varepsilon_{r2} = 16.00$).

To strictly isolate the impact of the constitutive contrast and rule out confounding effects of adaptive grid refinement ($h$-refinement), \textbf{the exact same point cloud distribution and background integration grid were employed across all four cases} ($N_{x1}=11$, $N_{x2}=11$, $N_y=21$, maintaining an invariant nodal density over the $[\pi \times \pi]\text{ m}$ domain).

Table~\ref{tab:contrast} summarizes the fundamental cutoff wavenumbers computed by 2D VNMM against exact analytical values for each dielectric material, complementing the parametric curve shown earlier in Figure~\ref{fig:jump}c.

\begin{table}[H]
\centering
\caption{Accuracy of the fundamental cutoff wavenumber across different dielectric contrast levels with a fixed point cloud.}
\label{tab:contrast}
\resizebox{\textwidth}{!}{
\begin{tabular}{lcccc}
\toprule
Dielectric Material & Permittivity $\varepsilon_{r2}$ & Exact Analytical $k_c$ & VNMM 2D $k_c$ & Relative Error (\%) \\
\midrule
\textbf{Teflon} & $2.25$ & $0.7709\text{ rad/m}$ & $0.7628\text{ rad/m}$ & $\mathbf{1.05\%}$ \\
\textbf{Standard Dielectric} & $4.00$ & $0.6082\text{ rad/m}$ & $0.6033\text{ rad/m}$ & $\mathbf{0.80\%}$ \\
\textbf{Alumina} & $9.80$ & $0.4029\text{ rad/m}$ & $0.4006\text{ rad/m}$ & $\mathbf{0.58\%}$ \\
\textbf{Microwave Ceramic} & $16.00$ & $0.3183\text{ rad/m}$ & $0.3167\text{ rad/m}$ & $\mathbf{0.52\%}$ \\
\bottomrule
\end{tabular}}
\end{table}

The results in Table~\ref{tab:contrast} and Figure~\ref{fig:jump}c show that the fundamental mode error remained strictly below $1.1\%$ over the entire contrast range, accurately tracking the theoretical curve. A notable physical trend is that the relative error of the fundamental cutoff wavenumber decreases monotonically as the dielectric contrast increases ($1.05\% \to 0.52\%$). In the fundamental mode, higher values of $\varepsilon_{r2}$ strongly confine the electromagnetic energy within the dielectric slab ($\Omega_2$). The planar interface increasingly behaves like a high-impedance virtual magnetic boundary with respect to the air region ($\Omega_1$), producing a steep yet smooth evanescent decay in the air and rendering the core field variation in $\Omega_2$ well-resolved by the local polynomial basis.

For higher-order resonant modes characterized by significantly shorter wavelengths inside dense media ($\lambda_d = \lambda_0 / \sqrt{\varepsilon_{r2}}$), standard computational practice suggests local nodal refinement scaling with $\sqrt{\varepsilon_{r2}}$. Nevertheless, the fact that VNMM accurately captures the fundamental mode with sub-percent error under an invariant, unrefined point cloud highlights the inherent robustness of the support partition strategy against severe constitutive jumps.

\section{Benchmarks in Anisotropic Media and Complex Geometries}

To demonstrate the applicability of VNMM to advanced electromagnetic problems, two reference literature benchmarks were examined \cite{ortiz2023,jiang2014}, followed by a challenging four-material junction.

\subsection{Homogeneous Anisotropic H-Shaped Ridge Waveguide (Reentrant Corners)}
The first benchmark investigates a non-convex domain with reentrant corners that induce electric field singularities, filled by a gyrotropic medium with a complex Hermitian permittivity tensor:
\begin{equation}
\mu_r = 1.0, \quad \overline{\overline{\varepsilon}}_r = \begin{bmatrix} 1 & -j \\ j & 4 \end{bmatrix}
\end{equation}
The cross-section occupies $x \in [-2.5, 2.5]\text{ m}$ and $y \in [-1.5, 1.5]\text{ m}$, with central reentrant notches narrowing the channel to $|y| \le 0.5\text{ m}$. Discretization uses alternating $\pm 45^\circ$ interior directors and tangential PEC directors.

Figure~\ref{fig:geom_h} illustrates the geometry, constitutive tensors, and reentrant corners of the H-shaped waveguide, while Table~\ref{tab:hguide} compares the first 5 eigenvalues computed by 2D VNMM against high-density edge FEM (427,492 edges), Mixed FEM \cite{jiang2014}, and EMM \cite{lima2017}.

\begin{figure}[H]
\centering
\includegraphics[width=0.85\textwidth]{Geometria_Guia_H_Anisotropico.png}
\caption{Geometry and constitutive tensors for the homogeneous anisotropic H-shaped ridge waveguide with reentrant corners (Thesis Benchmark 4.3.2).}
\label{fig:geom_h}
\end{figure}

\begin{table}[H]
\centering
\caption{Comparison of the first 5 eigenvalues in the H-shaped waveguide between 2D VNMM and high-density edge FEM (427,492 edges).}
\label{tab:hguide}
\resizebox{\textwidth}{!}{
\begin{tabular}{cccccc}
\toprule
Mode & Ref. FEM (427k edges) & Mixed FEM \cite{jiang2014} & EMM \cite{lima2017} (1,529 edges) & 2D VNMM (Nodal Cloud) & Relative Error (\%) \\
\midrule
\textbf{1} & $0.041212$ & $0.041206$ & $0.040628$ & $0.041915$ & $\mathbf{1.71\%}$ \\
\textbf{2} & $0.400641$ & $0.400552$ & $0.392058$ & $0.412739$ & $\mathbf{3.02\%}$ \\
\textbf{3} & $0.786745$ & $0.786704$ & $0.783254$ & $0.774186$ & $\mathbf{1.60\%}$ \\
\textbf{4} & $0.998524$ & $0.998524$ & $0.987213$ & $0.914897$ & $\mathbf{8.38\%}$ \\
\textbf{5} & $1.638579$ & $1.638369$ & $1.615863$ & $1.638612$ & $\mathbf{0.00\%}$ \\
\bottomrule
\end{tabular}}
\end{table}

The results in Table~\ref{tab:hguide} confirm that 2D VNMM computes the waveguide eigenvalues accurately without spurious modes, effectively overcoming the field singularities caused by sharp reentrant corners and conducting edges. Even in the presence of a complex Hermitian gyrotropic tensor, the method achieved an error of only $1.71\%$ on the dominant mode and exact agreement ($0.00\%$) on the 5th resonant mode relative to the 400k+ DOF FEM reference.

\subsection{Bimaterial Anisotropic Waveguide with Internal Circular PEC Cylinders}
The second benchmark simultaneously incorporates material discontinuity, distinct anisotropic tensors in each region, and internal curved conductor boundaries:
\begin{itemize}
  \item \textbf{Geometry:} Rectangular cross-section $[-2, 2] \times [-1, 1]\text{ m}$ divided at $x = 0$ into $\Omega_1 = [-2, 0] \times [-1, 1]\text{ m}$ and $\Omega_2 = [0, 2] \times [-1, 1]\text{ m}$. Each subregion encloses a circular PEC cylinder of radius $r = 1/3\text{ m}$ centered at $(-1, 0)$ and $(1, 0)$;
  \item \textbf{Medium 1 ($\Omega_1$):} $\mu_{r1} = 1.0, \quad \overline{\overline{\varepsilon}}_{r1} = \operatorname{diag}(1, 2)$ (biaxial anisotropy);
  \item \textbf{Medium 2 ($\Omega_2$):} $\mu_{r2} = 3.0, \quad \overline{\overline{\varepsilon}}_{r2} = \begin{bmatrix} 3 & -2j \\ 2j & 8 \end{bmatrix}$ (gyrotropic anisotropy with $\mu_r \neq 1$).
\end{itemize}

Figure~\ref{fig:geom_circ} depicts the geometric and constitutive setup of the bimaterial waveguide with internal circular PEC inclusions and planar interface at $x = 0$, while Table~\ref{tab:circguide} compares the first 5 eigenvalues obtained by 2D VNMM against reference FEM (340,363 edges), Mixed FEM \cite{jiang2014}, and EMM \cite{lima2017}.

\begin{figure}[H]
\centering
\includegraphics[width=0.85\textwidth]{Geometria_Guia_Bimaterial_Circ.png}
\caption{Geometry of the bimaterial anisotropic waveguide with internal cylindrical PEC inclusions and planar interface at $x = 0$ (Thesis Benchmark 4.3.3).}
\label{fig:geom_circ}
\end{figure}

\begin{table}[H]
\centering
\caption{Comparison of the first 5 eigenvalues in the bimaterial waveguide with PEC inclusions between 2D VNMM and reference FEM (340,363 edges).}
\label{tab:circguide}
\resizebox{\textwidth}{!}{
\begin{tabular}{cccccc}
\toprule
Mode & Ref. FEM (340k edges) & Mixed FEM \cite{jiang2014} & EMM \cite{lima2017} (8,566 edges) & 2D VNMM (Nodal Cloud) & Relative Error (\%) \\
\midrule
\textbf{1} & $0.051876$ & $0.051841$ & $0.051825$ & $0.051641$ & $\mathbf{0.45\%}$ \\
\textbf{2} & $0.205316$ & $0.205012$ & $0.204686$ & $0.203056$ & $\mathbf{1.10\%}$ \\
\textbf{3} & $0.247201$ & $0.247126$ & $0.246919$ & $0.255165$ & $\mathbf{3.22\%}$ \\
\textbf{4} & $0.423569$ & --- & $0.423277$ & $0.414830$ & $\mathbf{2.06\%}$ \\
\textbf{5} & $0.625603$ & --- & $0.624568$ & $0.637011$ & $\mathbf{1.82\%}$ \\
\bottomrule
\end{tabular}}
\end{table}

Table~\ref{tab:circguide} demonstrates that 2D VNMM finds all physical eigenvalues without spurious-mode contamination, despite the combined complexity of curved internal PEC boundaries, a discontinuous material interface at $x = 0$, and distinct anisotropic tensors. Aligning nodal directors tangentially along both the circular PEC boundaries and the planar interface enforces boundary conditions and tangential continuity between anisotropic media, achieving a dominant mode error of only $\mathbf{0.45\%}$ and errors below $3.3\%$ across all modes.

\subsection{Cavity with Four-Material Corner Junction (\textit{Quad-Junction})}
To stress-test the method against multiple orthogonal interfaces intersecting at a single point, a square PEC cavity $\Omega = [-1, 1] \times [-1, 1]\text{ m}$ divided into four distinct dielectric quadrants meeting at $(0, 0)$ was modeled ($\mu_r = 1.0$):
\begin{itemize}
  \item $Q_1$ ($x < 0, y > 0$): $\varepsilon_{r1} = 1.0$;
  \item $Q_2$ ($x > 0, y > 0$): $\varepsilon_{r2} = 4.0$;
  \item $Q_3$ ($x < 0, y < 0$): $\varepsilon_{r3} = 2.0$;
  \item $Q_4$ ($x > 0, y < 0$): $\varepsilon_{r4} = 9.0$ ($9:1$ maximum contrast at the central vertex).
\end{itemize}

The point cloud ($624$ total nodes: $528$ active, $96$ staggered PEC nodes, $3\times 3$ Gauss quadrature, $s_{\mathrm{div}} = 10.0$) uses tangential directors along the four semi-interfaces ($x = 0$ and $y = 0$), omitting a node exactly at $(0, 0)$ to prevent over-constraining the vertex. As an independent high-precision reference, the scalar $H_z$ Helmholtz formulation was solved using fine FEM ($Q_1$, $100 \times 100 = 10,201$ nodes), naturally free of spurious modes.

Figure~\ref{fig:geom_quad} shows the point cloud distribution, alternating $\pm 45^\circ$ interior directors, and tangential interface directors with strict quadrant-based support partitioning, while Table~\ref{tab:quad} compares the first 6 eigenvalues computed by 2D VNMM against the fine FEM $H_z$ reference.

\begin{figure}[H]
\centering
\includegraphics[width=0.72\textwidth]{Geometria_Nuvem_Quad_Junction.png}
\caption{Cavity with four-material cross junction meeting at $(0,0)$, showing point cloud distribution and tangential interface directors.}
\label{fig:geom_quad}
\end{figure}

\begin{table}[H]
\centering
\caption{Comparison of the first 6 eigenvalues in the four-material junction cavity between 2D VNMM and fine FEM $H_z$ ($100\times 100$).}
\label{tab:quad}
\begin{tabular}{ccccc}
\toprule
Mode & Ref. FEM ($H_z$, $100\times 100$) & 2D VNMM (Nodal Cloud) & Relative Error (\%) & $R_{\mathrm{div}}$ (\%) \\
\midrule
\textbf{1} & $0.515505$ & $0.513771$ & $\mathbf{0.34\%}$ & $0.54\%$ \\
\textbf{2} & $0.747157$ & $0.690568$ & $7.57\%$ & $0.71\%$ \\
\textbf{3} & $1.718475$ & $1.734055$ & $\mathbf{0.91\%}$ & $0.36\%$ \\
\textbf{4} & $2.026521$ & $2.003816$ & $\mathbf{1.12\%}$ & $1.22\%$ \\
\textbf{5} & $2.554525$ & $2.474523$ & $\mathbf{3.13\%}$ & $1.75\%$ \\
\textbf{6} & $3.788495$ & $3.922816$ & $\mathbf{3.55\%}$ & $1.64\%$ \\
\bottomrule
\end{tabular}
\end{table}

The results in Table~\ref{tab:quad} and Figure~\ref{fig:geom_quad} confirm that VNMM captures the multi-corner physics without spurious modes ($R_{\mathrm{div}} < 1.8\%$ across all modes): dominant mode error is only $\mathbf{0.34\%}$ and modes 3 and 4 exhibit errors around $\mathbf{1\%}$, proving that strict support partitioning and tangential interface directors generalize naturally to singular multi-material vertices with high contrast ($9:1$).

\section{Conclusions}

The development and systematic validation of the \textbf{2D Vector Nodal Meshless Method (VNMM)} on scattered point clouds confirms its capability for solving complex electromagnetic eigenvalue problems:
\begin{enumerate}
  \item \textbf{Reproducing Completeness and Nodal Invariance:} Directional vector projection guarantees first-order polynomial consistency ($\mathcal{P}^1$) to machine precision ($10^{-16}$) regardless of geometric point disorder (0\% to 50\% jitter). This occurs because the local collocation condition ($\mathbf{N}_i(\mathbf{x}_k)\cdot\mathbf{t}_k = \delta_{ik}$) is enforced exactly on a complete $\mathcal{P}^1$ vector basis at each evaluation point, making linear polynomial reproduction independent of Cartesian alignment.
  \item \textbf{Benefit of Irregular Point Clouds:} Moderately irregular point clouds (10\% to 25\% jitter) reduce the mean spectral error by more than 50\% (down to $\approx 1.1\%$) compared to rigid regular grids while keeping $R_{\mathrm{div}} < 3.3\%$. Physically, breaking regular Cartesian symmetry diversifies relative stencil distances and angles, disrupting directional numerical dispersion (\textit{grid-locking}), while divergence regularization ($s_{\mathrm{div}}\mathbf{K}_{\mathrm{div}}$) shifts irrotational gradient modes out of the physical spectrum.
  \item \textbf{Discontinuity and Corner Resolution without Smearing:} Strict material-based support partitioning prevents cross-material interpolation, while assigning a single tangential director ($\mathbf{t}_j = \boldsymbol{\tau}$) per interface node enforces exact $E_t$ continuity and leaves $E_n$ free to undergo the natural physical jump ($\varepsilon_1 E_{1n} = \varepsilon_2 E_{2n}$) without Gibbs oscillations---enabling $0.34\%$ accuracy even at singular four-material intersections with $9:1$ contrast.
  \item \textbf{Broad Applicability to Anisotropic Media, Singularities, and PEC Boundaries:} Conformity in $H(\mathbf{curl})$ combined with divergence penalization and staggered tangential PEC directors yields spurious-free eigenvalues in waveguides with complex Hermitian gyrotropic tensors, overcoming field singularities at reentrant corners ($1.71\%$ error in the H-shaped waveguide) and modeling curved internal PEC inclusions in anisotropic media ($0.45\%$ dominant mode error).
\end{enumerate}

\begin{thebibliography}{9}
\bibitem{harrington2001} R. F. Harrington, \textit{Time-Harmonic Electromagnetic Fields}. IEEE Press / Wiley-Interscience, 2001.
\bibitem{collin1991} R. E. Collin, \textit{Field Theory of Guided Waves}. IEEE Press, 1991.
\bibitem{jin2014} J. M. Jin, \textit{The Finite Element Method in Electromagnetics}. John Wiley \& Sons, 2014.
\bibitem{belytschko1994} T. Belytschko, Y. Y. Lu, and L. Gu, ``Element-free Galerkin methods,'' \textit{Int. J. Numer. Methods Eng.}, vol. 37, no. 2, pp. 229--256, 1994.
\bibitem{ortiz2023} L. A. Garcia Ortiz, \textit{M\'etodo Sem Malha Nodal Utilizando Fun\c{c}\~oes de Forma Vetoriais Baseadas em H(curl)}. Ph.D. Dissertation, Universidade Federal de Minas Gerais (UFMG), Belo Horizonte, Brazil, 2023.
\bibitem{ortiz2021} L. A. G. Ortiz, R. C. Mesquita, and N. Z. Lima, ``Nodal meshless method with vectorial shape functions based on H(curl),'' \textit{IEEE Trans. Magn.}, vol. 57, no. 6, pp. 1--4, 2021.
\bibitem{jiang2014} W. Jiang, N. Liu, Y. Tang, and Q. H. Liu, ``Mixed finite element method for 2D vector Maxwell's eigenvalue problem in anisotropic media,'' \textit{Prog. Electromagn. Res.}, vol. 148, pp. 159--170, 2014.
\bibitem{lima2017} N. Z. Lima and R. C. Mesquita, ``Edge meshless method applied to vector electromagnetic problems,'' \textit{IEEE Trans. Magn.}, vol. 53, no. 6, pp. 1--4, 2017.
\bibitem{mesquita_eabe} R. C. Mesquita, E. J. da Silva, B. M. F. Gon\c{c}alves, and L. A. G. Ortiz, ``A stabilized vector nodal meshless method for $H(\mathbf{curl})$ 2D electromagnetic problems on irregular point clouds,'' \textit{Eng. Anal. Bound. Elem.} (submitted), 2026.
\end{thebibliography}

\end{document}
